% DRAFT related-work section (Efe). Rewrite in the group's own words before
% submission -- see docs/framework.md: "the group writes the report in its own words".
\section{Related Work}
\label{sec:related}

\subsection{Dimensional music emotion recognition and DEAM}
Music emotion recognition (MER) either assigns discrete emotion classes or
places a piece in a continuous space; most recent work follows Russell's
circumplex model~\cite{russell1980}, in which valence (negative--positive) and
arousal (calm--energetic) span the plane. Early surveys already identified the
subjectivity of the labels and the small size of annotated corpora as the main
obstacles~\cite{kim2010}. Yang \emph{et al.}~\cite{yang2008} cast MER as a
regression problem on valence and arousal and trained support vector regression
(SVR) on hand-crafted audio descriptors, which became the standard baseline.
The MediaEval ``Emotion in Music'' tasks (2013--2015) produced DEAM, the largest
openly licensed dataset with both static (per-clip) and dynamic (2\,Hz) valence
and arousal ratings, together with frame-level openSMILE
features~\cite{aljanaki2017,eyben2010}. Its organisers report that arousal is
predicted far more reliably than valence, a gap that has persisted in later
work. G\'omez-Ca\~n\'on \emph{et al.}~\cite{gomezcanon2021} argue that MER
should move beyond a single ``ground truth'' and treat listener variability and
reproducible evaluation as first-class concerns, which motivates the
uncertainty part of this project.

\subsection{Hand-crafted features versus learned representations}
On DEAM, openSMILE descriptors with linear or kernel regressors remain
competitive; Cheuk \emph{et al.}~\cite{cheuk2020} report $R^2$ of about 0.32
(valence) and 0.64 (arousal) for such baselines and improve on them with a
triplet network. Choi \emph{et al.}~\cite{choi2017} showed that features of a
convolutional network pretrained on music tagging transfer to emotion
regression with only a shallow head. This idea has since scaled up.
PANNs~\cite{kong2020} are CNNs trained on the 2M-clip AudioSet that transfer to
many audio tasks. Castellon \emph{et al.}~\cite{castellon2021} probed the
generative model Jukebox with shallow heads and found that its representations
beat tagging-pretrained ones on emotion regression, among other music
information retrieval (MIR) tasks. MERT~\cite{li2024mert} is a self-supervised
music model trained with acoustic and pitch (CQT) teachers; in the MARBLE
benchmark~\cite{yuan2023marble} frozen MERT layers with a probing head are
among the strongest representations for emotion regression. Kang and
Herremans~\cite{kang2025} build on MERT features together with chord and key
information for unified dimensional and categorical MER. A common thread is
the \emph{frozen encoder + small head} protocol, which suits small datasets such
as DEAM and isolates the effect of the representation. A limitation is that
most of these studies use random (not artist-disjoint) splits, so reported
gains may partly reflect artist or production cues; we therefore evaluate on
artist-grouped folds.

\subsection{Heteroscedastic regression}
Nix and Weigend~\cite{nix1994} proposed networks that output both a mean and a
variance and are trained with the Gaussian negative log-likelihood (NLL).
Kendall and Gal~\cite{kendall2017} distinguished the input-dependent noise such
a head captures (aleatoric) from model uncertainty (epistemic), which a single
network does not represent; deep ensembles~\cite{lakshminarayanan2017} estimate
the latter through the spread of independently trained members. Training with
NLL is known to be fragile: the variance term can down-weight hard examples and
let the mean underfit. Seitzer \emph{et al.}~\cite{seitzer2022} propose the
$\beta$-NLL loss, which re-weights each sample by $\sigma^{2\beta}$, and Stirn
\emph{et al.}~\cite{stirn2023} propose a ``faithful'' variant that keeps the mean
fit identical to an MSE model. For evaluation, Levi \emph{et
al.}~\cite{levi2022} show that average calibration can hide poor per-sample
behaviour and recommend binned reliability diagrams (ENCE), which we adopt.

\subsection{Uncertainty and annotator disagreement in MER}
The closest work to ours is Watcharasupat \emph{et al.}~\cite{watcharasupat2025},
who benchmark probabilistic losses, Monte-Carlo dropout and ensembles for MER
and find that models learn the central tendency of the ratings well but
largely fail to reproduce the empirical spread of listener responses. Their
study does not explain \emph{why}. Our project positions itself as a partial
replication and extension: (i) we make explicit that a variance learned against
the \emph{mean} rating estimates the model's error on that mean, not the spread
of individual listeners -- the distinction raised in the instructor's feedback;
and (ii) we measure how reliably the per-clip listener spread can be estimated
at all in DEAM before asking a model to predict it.

\subsection{Evaluation practice}
Bouthillier \emph{et al.}~\cite{bouthillier2021} show that benchmark
comparisons are dominated by data-split and seed variance, and recommend
comparing models over many paired splits with a probability-of-improvement
criterion rather than single-split point estimates. This motivates our
switch from one 70/15/15 split to artist-grouped 10-fold cross-validation with
paired per-fold comparisons.
